\documentclass[11pt, a4paper]{article}
\usepackage{amsmath, amssymb, booktabs, geometry, hyperref, xcolor, listings}
\geometry{margin=2.5cm}
\hypersetup{colorlinks=true, linkcolor=blue}
\lstset{basicstyle=\ttfamily\small, backgroundcolor=\color{gray!10}, frame=single}

\title{\textbf{EGARCH(1,1,1) Model}\\[4pt]
\large Portfolio VaR Engine --- Model Documentation}
\author{Risk Modelling --- Trade Republic}
\date{April 2026}

\begin{document}
\maketitle

%-----------------------------------------------------------------------
\section{Overview}

The Exponential GARCH model (EGARCH; Nelson, 1991) models the
\textit{logarithm} of the conditional variance, rather than the level.
This has two key advantages over GJR-GARCH: (1) positivity of $\sigma_t^2$
is guaranteed without parameter constraints, and (2) both the magnitude
and the sign of past shocks affect future volatility, naturally capturing
the leverage effect.

\textbf{Engine key:} \texttt{'egarch'} \\
\textbf{Sources:}
\begin{itemize}
  \item \texttt{risk\_models/garch\_models.py}, lines 16--35 (\texttt{garch\_exp})
  \item \texttt{risk\_models/garch\_models.py}, lines 58--131 (\texttt{calculate\_vaR\_garch})
  \item \texttt{risk\_analytics/risk\_engines.py}, lines 135--152 (engine dispatch)
\end{itemize}
\textbf{Library:} \texttt{arch} (Kevin Sheppard)

%-----------------------------------------------------------------------
\section{Volatility Model}

\subsection{EGARCH(1,1,1) Specification}

Let $r_t$ denote the portfolio return at time $t$. The model assumes:
\begin{equation}
r_t = \mu + \varepsilon_t, \qquad
\varepsilon_t = \sigma_t z_t, \qquad
z_t \overset{i.i.d.}{\sim} D(0,1)
\end{equation}
The log-variance dynamics are:
\begin{equation}
\log \sigma_t^2
= \omega
+ \alpha \left(\frac{|\varepsilon_{t-1}|}{\sigma_{t-1}}
              - \mathbb{E}\!\left[\frac{|\varepsilon_{t-1}|}{\sigma_{t-1}}\right]\right)
+ \gamma \frac{\varepsilon_{t-1}}{\sigma_{t-1}}
+ \beta \log \sigma_{t-1}^2
\label{eq:egarch}
\end{equation}

\begin{itemize}
  \item $\omega$: intercept of the log-variance equation
  \item $\alpha \geq 0$: magnitude effect (symmetric response to shock size)
  \item $\gamma$: sign (asymmetry/leverage) effect;
        $\gamma < 0$ implies negative shocks increase volatility more than
        positive shocks
  \item $\beta$: persistence in log-variance; $|\beta| < 1$ ensures stationarity
  \item $\mathbb{E}[|z_{t-1}|]$ is a normalisation constant that depends on $D$
\end{itemize}

Because the variance is modelled in logarithms, \textit{no non-negativity
constraints} are needed on $\omega$, $\alpha$, $\gamma$, or $\beta$.

\subsection{Comparison with GJR-GARCH}

\begin{tabular}{lll}
\toprule
Feature & GJR-GARCH & EGARCH \\
\midrule
Variance modelled & Level $\sigma_t^2$ & Log $\log\sigma_t^2$ \\
Parameter constraints & Yes ($\omega, \alpha, \beta \geq 0$) & None \\
Leverage effect & Indicator $\mathbb{1}_{\varepsilon<0}$ & Sign term $\gamma z_{t-1}$ \\
Stationarity condition & $\alpha + \frac{1}{2}\gamma + \beta < 1$ & $|\beta| < 1$ \\
\bottomrule
\end{tabular}

\subsection{ARCH library call}

\begin{lstlisting}[language=Python]
arch_model(rets_rescaled, p=1, q=1, o=1, vol='EGARCH', dist=dist)
\end{lstlisting}
Parameters \texttt{p=1}, \texttt{q=1} set the ARCH and GARCH lag orders;
\texttt{o=1} activates the asymmetric (leverage) term.

%-----------------------------------------------------------------------
\section{VaR and CVaR Calculation}

The VaR and CVaR computation is \textbf{identical} to the GJR-GARCH engine,
both calling \texttt{calculate\_vaR\_garch()} with \texttt{garch\_engine='egarch'}.
For completeness, the formulas are reproduced below.

\subsection{One-Step Forecast}

\begin{equation}
\hat{\mu}_{t+1} = \hat{\mu}, \qquad \hat{\sigma}_{t+1}^2 = \text{forecast variance}
\end{equation}

\subsection{Parametric VaR (without FHS)}

\textbf{Value-at-Risk:}
\begin{equation}
\widehat{\text{VaR}}_\alpha
= -\hat{\mu}_{t+1} + \hat{\sigma}_{t+1} \cdot q_\alpha^D
\end{equation}

\textbf{CVaR --- Normal:}
\begin{equation}
\widehat{\text{CVaR}}_\alpha^{\text{N}}
= -\hat{\mu}_{t+1} + \hat{\sigma}_{t+1} \cdot
  \frac{\phi\!\left(\Phi^{-1}(\alpha)\right)}{\alpha}
\end{equation}

\textbf{CVaR --- Student-$t$:}
\begin{equation}
\widehat{\text{CVaR}}_\alpha^{t}
= -\hat{\mu}_{t+1}
  - \frac{\hat{\sigma}_{t+1}}{\alpha(1 - \hat{\nu})}
    \Bigl(\hat{\nu} - 2 + \bigl[t^{-1}(\alpha;\hat{\nu})\bigr]^2\Bigr)
    \cdot f_t\!\bigl(t^{-1}(\alpha;\hat{\nu});\hat{\nu}\bigr)
\end{equation}

\subsection{Filtered Historical Simulation (FHS)}

When \texttt{fhs=True}, standardised residuals
$\hat{z}_t = (r_t - \hat{\mu})/\hat{\sigma}_t$ replace the parametric
distribution quantile:
\begin{equation}
q_\alpha^{\text{FHS}} = Q_{1-\alpha}\!\left(\{\hat{z}_t\}\right)
\end{equation}

\subsection{Holding-Period Scaling}
\begin{equation}
\widehat{\text{VaR}}_\alpha^{(H)} = \widehat{\text{VaR}}_\alpha^{(1)} \cdot \sqrt{H}
\end{equation}

%-----------------------------------------------------------------------
\section{Parameters}

\begin{tabular}{llll}
\toprule
Parameter & Type & Default & Description \\
\midrule
\texttt{window}   & int   & 250   & Lookback window \\
\texttt{rescale}  & float & 1000  & Rescaling factor for numerical stability \\
\texttt{decay}    & float & 1     & EWMA pre-weighting decay ($\lambda=1$: no decay) \\
\texttt{ftol}     & float & 1e-6  & Optimizer convergence tolerance \\
\texttt{max\_iter}& int   & 150   & Maximum optimizer iterations \\
\texttt{fhs}      & bool  & False & Use Filtered Historical Simulation \\
\texttt{distr}    & str   & \texttt{'Normal'} & Innovation distribution \\
\texttt{qtls}     & list  & [0.05, 0.01, 0.001] & Tail probability levels \\
\texttt{holding\_period} & int & 1 & Forecast horizon in days \\
\bottomrule
\end{tabular}

%-----------------------------------------------------------------------
\section{Advantages and Limitations}

\begin{tabular}{p{6cm} p{8cm}}
\toprule
Advantages & Limitations \\
\midrule
No non-negativity constraints on parameters & More complex likelihood than GJR \\
Natural leverage effect via sign term & Harder to interpret parameters \\
Log-variance avoids negative variance & Sensitive to initialisation \\
Well-suited to equity return series & Square-root scaling is approximate \\
\bottomrule
\end{tabular}

%-----------------------------------------------------------------------
\section{Usage Example}

\begin{lstlisting}[language=Python]
from risk_analytics import risk_engines as re

var_egarch = re.portfolio_vaR(
    rets_orig  = rets,
    wgts       = weights,
    engine     = 'egarch',
    fmt_engine = {
        'window'         : 250,
        'rescale'        : 1000,
        'decay'          : 1,
        'ftol'           : 1e-6,
        'max_iter'       : 150,
        'fhs'            : False,
        'distr'          : 't',
        'qtls'           : [0.05, 0.01, 0.001],
        'holding_period' : 1,
    }
)
\end{lstlisting}

%-----------------------------------------------------------------------
\section{References}

Nelson, D.B. (1991).
\textit{Conditional Heteroskedasticity in Asset Returns: A New Approach.}
Econometrica, 59(2), 347--370.

\end{document}
